\section{Method}
\label{sec:method}

\noindent\textbf{Input.} A target volume $I^t \in \mathbb{R}^{D\times H\times W}$
and a context set $\{(I^c_k, L^c_k)\}_{k=1}^{K}$ of image/mask pairs for
the same anatomical class as the (unknown) target mask.

\noindent\textbf{Output.} A predicted target mask $\hat{L}^t$, produced
by a coarse-to-fine cascade over $M$ resolution levels
$\{r_1, \dots, r_M\}$. \cref{fig:overview} contrasts our architecture
with Medverse's.

\begin{figure*}[t]
  \centering
  \begin{subfigure}[b]{0.48\linewidth}
    \centering
    \includegraphics[width=\linewidth]{arch_medverse.pdf}
    \caption{Medverse: early concatenation + dual U-Net.}
    \label{fig:overview-medverse}
  \end{subfigure}
  \hfill
  \begin{subfigure}[b]{0.48\linewidth}
    \centering
    \includegraphics[width=\linewidth]{arch_patchset.pdf}
    \caption{Ours: separate image/label streams + bi-axial attention.}
    \label{fig:overview-ours}
  \end{subfigure}
  \caption{Architecture overview, contrasted against Medverse's
  early-concatenation fusion. \TODO{draft diagram; finalize labels/scale
  once the architecture is settled.}}
  \label{fig:overview}
\end{figure*}

\subsection{Backbone}
\label{sec:backbone}

Each volume (target or context) is tokenized by a shared encoder --
either the nnU-Net residual-encoder stages
(\texttt{ResEncInContext3D})~\cite{isensee2021nnu} or a 3D patch
embedding (\texttt{ViTInContext3D}) -- into a grid of spatial tokens
with 3D axial rotary position embeddings (RoPE), following the
architecture already implemented in our codebase. At the bottleneck, a
two-stage attention block first lets tokens attend within their own
volume (target and each context volume independently, sharing weights),
then lets the target cross-attend, read-only, into all $K$ context
bottlenecks. \TODO{restate/adjust to match final architecture diagram.}

\subsection{Bi-Axial Image--Label Attention}
\label{sec:biaxial}

Rather than folding the label into the image as a second input channel
before encoding (concatenation, as in Medverse) or fusing image and
label features once via a pixel shuffle-unshuffle step (Iris), we keep
image tokens $\mathbf{z}^{\text{img}}$ and label tokens
$\mathbf{z}^{\text{lbl}}$ as separate streams and attend over two axes:

\begin{itemize}[leftmargin=1.4em,itemsep=0.1em]
  \item \textbf{Patch axis}: standard spatial self-/cross-attention
    among tokens at the same image/label identity, as described in
    \cref{sec:backbone}.
  \item \textbf{Image/label axis}: at each spatial location, the image
    and label tokens attend to each other, letting label evidence
    sharpen the corresponding image features and vice versa.
\end{itemize}

\cref{fig:biaxial} illustrates both axes within a single layer.

The target's label stream starts from a query prior (\cref{sec:cascade})
rather than a zero/blank mask, so this axis is informative for the
target as well as for the contexts. \TODO{formalize with an equation once
the exact factorization (interleaved vs.\ parallel bi-axial layers) is
finalized; specify how many layers alternate between the two axes.}

\begin{figure}[t]
  \centering
  \includegraphics[width=0.9\linewidth]{arch_biaxial.pdf}
  \caption{Bi-axial attention within one transformer layer, shown for a
  single context volume ($K{=}1$ here for clarity; in general $K$
  context volumes are each processed this way) and 3 of its $N$ spatial
  tokens per stream. Left: the \emph{image/label axis} -- image and
  label tokens attend to each other within the same cell. Right: the
  \emph{patch axis} -- tokens attend across cells and, for the target,
  across contexts (edges stand in for full attention among all $N$
  tokens of both volumes). Gray edges are context-only; green edges
  touch the target row, which is architecturally distinct -- it
  cross-attends read-only into context, and its label-stream input is
  the predicted query prior (\cref{sec:cascade}) rather than ground
  truth. \TODO{draft diagram; finalize once the layer factorization
  above is settled.}}
  \label{fig:biaxial}
\end{figure}

\subsection{Coarse-to-Fine Cascade}
\label{sec:cascade}

Our model processes the target through $M$ increasingly fine
resolution levels $\{r_1 < \dots < r_M\}$ (e.g. a coarse full-volume pass
followed by finer crops). Two mechanisms connect consecutive levels.

\vspace{0.3em}\noindent\textbf{Query image prior.}
At level $\ell > 1$, the target's label tokens are initialized from the
upsampled prediction of level $\ell - 1$,
\begin{equation}
  \mathbf{z}^{\text{lbl}}_{\ell,0} = \operatorname{Embed}\bigl(\operatorname{Upsample}(\hat{L}^t_{\ell-1})\bigr),
  \label{eq:query-prior}
\end{equation}
instead of entering the network through a separate autoregressive branch
(Medverse) or not being used at all (Iris). This makes the previous
level's belief part of the same token stream that attends bi-axially
with the image, rather than a side channel fused later. \cref{fig:query-prior}
shows this flow across two levels.

\begin{figure*}[t]
  \centering
  \includegraphics[width=\linewidth]{arch_query_prior.pdf}
  \caption{Query image prior across two cascade levels, with the full
  model input/output flow. At level 1 the target has no real prior yet
  (dashed support-mean placeholder); at level 2 the target's second lane
  is a real crop of level 1's own prediction $\hat{L}^t_1$ -- the same
  lane the context mask always occupies. \TODO{draft diagram; finalize
  once the number of levels $M$ is settled.}}
  \label{fig:query-prior}
\end{figure*}

\vspace{0.3em}\noindent\textbf{Region-restricted fine processing.}
Given level $\ell-1$'s prediction, we take its predicted foreground
region (thresholded and padded by a margin) and restrict the sliding
window at level $\ell$ to that region, rather than re-scanning the full
volume as in Medverse's autoregressive steps. This lets compute at fine
levels scale with the size of the structure of interest rather than the
size of the input volume, at the cost of depending on the coarse level's
localization being approximately correct. \TODO{specify margin/threshold
values and failure-mode handling (e.g. what happens when the coarse
level misses the structure entirely).}

\subsection{Cross-Level Training via Register Tokens}
\label{sec:registers}

To connect the cascade with a gradient path that crosses resolution
levels -- rather than levels being independently supervised and only
chained via the rendered prediction in \cref{eq:query-prior} -- we add a
small set of learned register tokens $\mathbf{r}_\ell \in \mathbb{R}^{N_r \times d}$
at each level. Registers participate in the bi-axial attention at level
$\ell$ alongside image and label tokens, and their updated state is
carried forward, via a learned projection and a level-type embedding,
into level $\ell+1$'s register initialization:
\begin{equation}
  \mathbf{r}_{\ell+1,0} = W_r\, \mathbf{r}_{\ell,\text{final}} + \mathbf{e}_{\ell+1}.
  \label{eq:register-carry}
\end{equation}
This allows the model to propagate information that is useful for
finer-level processing but not fully captured by the rendered
segmentation mask alone (e.g. uncertainty or contextual cues about
neighboring structures). \cref{sec:cross-level-ablation} ablates this
mechanism against training levels independently with the query prior as
the only cross-level connection.

\subsection{Training Objective}
\label{sec:objective}

Each level is supervised at its own resolution with a combined
binary cross-entropy and soft Dice loss,
\begin{equation}
  \mathcal{L}_\ell = \mathcal{L}_{\text{BCE}}^\ell + \mathcal{L}_{\text{Dice}}^\ell,
\end{equation}
and the total loss sums over levels, \TODO{state whether uniform or
level-weighted, matching the final training config}:
\begin{equation}
  \mathcal{L} = \sum_{\ell=1}^{M} \lambda_\ell\, \mathcal{L}_\ell.
\end{equation}

\subsection{Synthetic Supervision}
\label{sec:synth}

A fraction $p_{\text{synth}}$ of training tasks are drawn from a
GMM-based synthetic generator instead of real (image, mask) pairs: a
real-anatomy supervoxel cohort is repainted with per-class Gaussian
intensities, giving a label-perfect (image, mask) pair with a novel
appearance. Paint noise is i.i.d.\ per voxel by default; an optional
multi-octave correlated variant (coarse-to-fine random fields, summed)
better matches the spatial autocorrelation of real tissue (measured
$\approx 0.6$ in real CT/MRI vs.\ $\approx 0.0$ for i.i.d.\ paint at
native resolution) without changing the per-class intensity variance.
\TODO{figure of a synth\_gmm training pair (i.i.d.\ vs.\ multi-octave
texture noise) next to a real (image, mask) pair; equation for the paint
model once finalized.}
